
\subsubsection{6.1 Interpretation of Findings}

The verbosity coefficient reversal (Δβ_L = +0.080, non-overlapping 95\% CI) is the most robust finding of this study. It is replicable without annotator metadata — requiring only preference labels and response text — and survives the Q_early stability check (β_Q = −0.017). The pattern is directionally consistent with the hypothesis that annotators in later phases of multi-round RLHF annotation assign greater positive weight to verbosity, which is a salient stylistic feature of AI-generated text.

The between-group tercile analysis on WebGPT provides supporting evidence at the population level: AI-typicality projection scores differ significantly across groups (F = 33.226), and this difference is specific to the AI-typicality direction (discriminant validity confirmed). However, the between-group design cannot distinguish individual annotator adaptation from group composition differences, and the pre-registered within-annotator effect-size threshold was not met.

The directional consistency across all three stylistic features (sign_consistent = true for β_L, β_H, β_S) is informative as a diagnostic: all deltas are positive despite only one meeting the CI non-overlap criterion. Hedging and structured reasoning may also exhibit drift but at magnitudes below detection given the available proxy design and topic confound.

\subsubsection{6.2 Limitations}

\textbf{L1: Index-based round stratification.} HH-RLHF does not contain per-annotation timestamps or annotator identifiers. Round stratification is based on row index, not genuine temporal ordering. The observed verbosity coefficient reversal may reflect between-cohort composition differences rather than within-annotator adaptation over time. This limits the causal interpretation to population-level association.

\textbf{L2: WebGPT between-group design.} The planned within-annotator PanelOLS regression was not executable due to the absence of worker identity fields in the public WebGPT release. The tercile proxy (|score_0 − score_1| as confidence proxy) collapsed to two effective groups, further limiting between-group power. Selection effects cannot be ruled out.

\textbf{L3: Partial AAI validation.} The AAI as specified comprises three components; only the first two are evaluated in this work. The reward-model behavioral divergence component (H-M3) and the downstream benchmark degradation test (H-M4) were not executed. Whether the observed stylistic coefficient shift propagates to reward model scoring behavior remains an open empirical question.

\textbf{L4: Topic distribution imbalance.} The chi-square test across HH-RLHF strata yields p = 4×10⁻²⁷⁵, indicating that strata differ substantially in topic distribution. This confound likely affects β_H and β_S more than β_L, since response length is less topic-specific than hedging or structural features.

\textbf{L5: H-M1 gate interpretation.} The MUST_WORK gate for H-M1 was evaluated on code executability, not scientific effect size. The pre-registered scientific criterion (effect size ≥ 0.1 SD per 1,000 tokens) was not met. Results should be interpreted as proof-of-concept for the measurement pipeline, not as validation of the full AAI component-2 specification.

\subsubsection{6.3 Implications}

If the verbosity coefficient reversal reflects genuine annotator adaptation rather than cohort composition effects, it implies that reward models trained on later annotation rounds optimize for a different preference target than those trained earlier. This temporal non-stationarity would be invisible to standard pairwise accuracy metrics.

The AAI provides a low-cost monitoring instrument for this phenomenon: round-stratified coefficient comparison requires only preference labels and response text. Practitioners with access to multi-round annotation datasets could apply this check as a quality gate prior to reward model training.

The null results from this study provide concrete, actionable guidance for future dataset design: confirming individual annotator causal adaptation requires genuine annotation timestamps, within-annotator session tracking across at least three sessions per annotator, and per-prompt disagreement labels to enable ambiguity stratification.

